\subsection{The descent operator on the quadratic annihilator}

Let $D=D_P:\GTz\longrightarrow\GTz$ be the continuous
homomorphism with the literal marking of (B.C1).  We use the
following discrete $\mathbb F_2$-vector spaces and linear maps:
\[
\begin{gathered}
 V=\{\eta\in\operatorname{Hom}_{\rm cts}(\GTz,\mathbb F_2):
                         \eta|_{\GT}=0\},\\
 D^*:V\longrightarrow V,\qquad D^*\eta=\eta\circ D,\\
 R_C:V\longrightarrow
       \operatorname{Hom}_{\rm cts}(C_P,\mathbb F_2),
       \qquad R_C(\eta)=\eta|_{C_P}.
\end{gathered}                                                   \tag{B.V1}
\]
The fixed subgroup $C_P=\{g:D(g)=g\}$ is closed.  The map $D^*$
has the indicated target by B.1.8: $D(F)=cFc$ for every $F\in\GT$.

\begin{result}{B.2.5}{Proved}{Descent and the quadratic annihilator}
The maps in (B.V1) satisfy
\[
 D^*V\cap\mathbb F_2\Delta=0,\qquad
 [\Delta]\ne0\text{ in }V/D^*V,\qquad R_C D^*=R_C.       \tag{B.V2}
\]
In particular, for $\chi_n=\Delta\circ D^n$ and every $n\geq1$,
\[
 \chi_n\in V,\qquad
 \mathbb F_2\Delta\cap
       \operatorname{span}_{\mathbb F_2}\{\chi_n:n\geq1\}=0,
 \qquad \chi_n|_{C_P}=\Delta|_{C_P}.                      \tag{B.V3}
\]
The nonzero character $\Delta+\chi_1$ vanishes on both $\GT$
and $C_P$.
\end{result}
\begin{proof}
By B.1.3 there is an actual element $k\in\ker D$ with
$\Delta(k)=1$.  Evaluation at this element is a well-defined
linear map $e_k:V\to\mathbb F_2$.  Every $\eta\in V$ gives
\[
                   e_k(D^*\eta)=\eta(D(k))=\eta(1)=0,
                   \qquad e_k(\Delta)=1.                \tag{B.V4}
\]
Thus the intersection and the quotient assertion in (B.V2)
follow by evaluating a putative equality at this same full
profinite element.  For $g\in C_P$ we have $D(g)=g$, so
$R_C(D^*\eta)(g)=\eta(g)=R_C(\eta)(g)$; this proves the third
identity as an equality of actual restriction maps.

B.1.8 and A.7 put every $\chi_n$ in $V$.  For $n\geq1$,
$\chi_n=D^*(\Delta\circ D^{n-1})$, so their span lies in
$D^*V$ and has zero intersection with $\mathbb F_2\Delta$.
Iteration of the restriction identity proves the last part
of (B.V3).  Finally $\Delta+\chi_1$ restricts to zero on
$C_P$ and on $\GT$, and its value at $k$ is one.
\end{proof}

B.2.8 below proves that $D^*\Delta$ is nonzero by constructing
an actual full profinite element on which it has value one.
Consequently $\Delta+\chi_1$ and $\Delta$ are distinct, and
$\Delta,\chi_1$ are linearly independent. B.2.24--B.2.25
subsequently prove that $D$ is surjective and that every
$\chi_n$ is independent of all the others. They determine
the full vector-space dimension and construct an actual
binary shift quotient. B.2.27 proves that the particular
family generated by $\Delta$ does not span $V$, through
an actual full descent-kernel element detected by
$\overline d_2$. B.2.28 constructs the two independent
shift families generated by these characters. Spanning
by that larger specified polynomial submodule is not
asserted.
